Тема диссертации находится на пересечении трёх направлений: оценки
долгосрочного эффекта онлайн-экспериментов, моделирования поведения
пользователей и методологии А/Б-тестирования. Ниже разбираются основные из
этих направлений, причём внимание уделяется не только их достижениям, но и
ограничениям, которые проявляются в задаче масштабируемой и интерпретируемой
оценки LTV.

\subsection*{Краткосрочные А/Б-эксперименты и прокси-метрики}

А/Б-тестирование давно закрепилось как стандартный инструмент оценки
продуктовых изменений на онлайн-платформах~\cite{kohavi2020trustworthy}.
Из-за операционных ограничений большинство экспериментов планируется так,
чтобы давать ответ за дни или недели, и команды вынуждены опираться на
краткосрочные прокси-метрики --- кликабельность (CTR), длительность сессии,
немедленную конверсию. Такие показатели удобно измерять, но они нередко
слабо связаны с долгосрочными бизнес-исходами.

Эта проблема систематически разбирается в литературе об ошибках
интерпретации онлайн-экспериментов. В работе~\cite{dmitriev2016dirty}
описана характерная подборка типичных <<ловушек>> (так называемая
<<dirty dozen>>), среди которых --- интерпретация краткосрочного выигрыша
как устойчивого долгосрочного эффекта. Хрестоматийный пример ---
увеличение объёма рекламных показов: оно даёт немедленный прирост выручки,
но со временем снижает удовлетворённость пользователей и повышает отток.
Именно такое рассогласование заставляет искать методы, способные переводить
краткосрочные сигналы в надёжные оценки долгосрочного эффекта.

\subsection*{Длительные эксперименты}

Самая очевидная стратегия --- продлить сам эксперимент. Опыт
Netflix~\cite{dimakopoulou2023evaluating} и
Microsoft~\cite{sadeghi2021novelty} показывает, что длительные дизайны
действительно позволяют наблюдать отложенные исходы, в первую очередь ---
изменения удержания и частоты покупок. Однако такие эксперименты дороги,
медленно дают ответ и оставляют бизнес под продолжительным риском.
К этому добавляется проблема \textit{выживаемости} (survivorship bias):
со временем в окне анализа остаются преимущественно наиболее вовлечённые
пользователи, что искажает оценку эффекта на уровне всей популяции.

\subsection*{Суррогатное моделирование долгосрочных эффектов}

Чтобы совместить быстрые циклы экспериментирования с потребностью в
долгосрочной оценке, развиваются методы \textit{суррогатного моделирования}:
долгосрочный исход прогнозируется по краткосрочным поведенческим сигналам.
По сути задача сводится к выбору или построению такой метрики (или
композитной метрики) на ранних этапах жизненного цикла пользователя, которая
сильно коррелирует с будущей ценностью. В индустриальной практике эта
конструкция получила название \textit{Surrogate
Index}~\cite{dimakopoulou2023evaluating}: она ускоряет принятие решений,
не разрывая их со стратегическими целями бизнеса.

Современные реализации опираются на машинное обучение и выучивают сложные
нелинейные отображения из ранних поведенческих признаков
в LTV~\cite{harper2024long}. На платформах Netflix и Microsoft такие модели
используются для прогнозирования последующего вовлечения и выручки по
ранним паттернам взаимодействия, а их калибровка проводится на исторических
А/Б-тестах для обеспечения
надёжности~\cite{dimakopoulou2023evaluating,sadeghi2021novelty}.

При всей выразительности эти модели обычно слабо интерпретируемы и требуют
значительного объёма исторических данных для калибровки. Кроме того, они
опираются на так называемое \textit{предположение суррогатности}:
краткосрочная прокси-метрика должна не только коррелировать с долгосрочным
исходом, но и полностью опосредовать причинный эффект воздействия на
него~\cite{dimakopoulou2023evaluating}. Формально это требование
сформулировано Prentice~\cite{prentice1989} в контексте клинических
испытаний: все причинные пути от воздействия к итоговому исходу должны
проходить исключительно через суррогат. Практика показывает, что
суррогатные индексы, обученные на краткосрочных метриках, успешно
предсказывают долгосрочные исходы --- это подтверждено валидацией более
чем на 200 А/Б-тестах~\cite{dimakopoulou2023evaluating}, --- однако само
предположение суррогатности трудно проверять, и оно может нарушаться из-за
эффектов новизны, привыкания и скрытых искажающих факторов.

\subsection*{Когортный анализ и поведенческая сегментация}

Когортный анализ относится к базовым техникам клиентской аналитики и
традиционно применяется для разложения LTV на компоненты: удержание,
частоту и денежную ценность (RFM)~\cite{rosset2003modeling}. Группируя
пользователей по общим признакам --- например, по месяцу регистрации или
уровню активности, --- команды получают возможность отслеживать эволюцию
ключевых метрик во времени.

В диссертации эта традиция продолжается, но акцент смещается с описательной
сегментации на \textit{целеориентированную}. В отличие от статических
когортных определений, орбитали выводятся из предсказательной задачи ---
прогнозирования будущей активности пользователя --- с помощью решающих
правил, извлечённых из ансамблевых древесных моделей машинного обучения,
широко применяемых на практике. Благодаря этому получаемые сегменты
остаются интерпретируемыми и одновременно настроены под задачу оценки
долгосрочного эффекта.

Важной теоретической опорой для связи RFM-признаков с пожизненной ценностью
служит работа Fader, Hardie и Lee~\cite{fader2005rfm}, в которой вводится
понятие \textit{iso-value curves} --- линий уровня в пространстве RFM,
объединяющих пользователей с одинаковой ожидаемой LTV. Авторы показывают,
что разные комбинации давности, частоты и денежной ценности могут
соответствовать сопоставимой долгосрочной ценности, и тем самым ставят под
сомнение упрощённые сегментационные схемы. Этот результат поддерживает
выбранный подход: вместо фиксированных RFM-эвристик отображение
поведенческих паттернов в LTV выучивается из данных, что обеспечивает более
точное и многогранное представление архетипов пользователей.

\subsection*{Параметрические модели поведения пользователей}

Существенный пласт литературы посвящён вероятностным моделям активности и
оттока. Среди наиболее известных --- Pareto/NBD и
BG/NBD~\cite{fader2005counting}, описывающие два сопряжённых процесса:
интенсивность транзакций и выбытие пользователя. Эти модели позволяют
прогнозировать будущие покупки и широко применяются в анализе клиентской
базы.

В рамках предварительного исследования мы применяли Pareto/NBD для описания
активности пользователей и оценки LTV, выраженной в днях вовлечённости.
Модель удовлетворительно воспроизводит агрегированные паттерны удержания,
но оказывается менее эффективной при работе с гетерогенной пользовательской
популяцией и короткими окнами наблюдения. К этому добавляется характерное
ограничение её предпосылок --- в частности, постоянной интенсивности
выбытия, --- которое нередко нарушается в динамичных продуктовых средах.

В более поздних работах изучается применение моделей времени до события и
анализа выживаемости для оценки долгосрочных
эффектов~\cite{luck2019learning}. Эти подходы, однако, требуют сложного
отбора признаков и плохо переносятся на широкий спектр типов экспериментов.

\subsection*{Uplift-моделирование и переходы между состояниями}

Uplift-моделирование нацелено на оценку индивидуального эффекта воздействия
и служит основой для персонализированных решений. Среди современных
подходов выделяется \textit{Uplift Random Forest}~\cite{guelman2015uplift},
расширяющий алгоритм случайного леса за счёт критерия расщепления,
ориентированного на гетерогенный эффект. Это позволяет выявлять подгруппы,
в которых воздействие оказывает наиболее сильное положительное или
отрицательное влияние, при сохранении устойчивости за счёт ансамблирования.
Несмотря на потенциал, такие методы чувствительны к шуму и требуют больших
объёмов выборки для надёжной оценки эффектов на индивидуальном уровне.

Предлагаемый подход дополняет uplift-моделирование, перенося фокус с
индивидуальных эффектов на \textit{групповые переходы} между поведенческими
состояниями. Описывая то, как воздействие сдвигает распределение
пользователей по орбиталям, мы получаем причинную интерпретацию,
устойчивую к шуму и применимую на практике. По духу подход близок к
марковским моделям отношений с клиентом~\cite{pfeifer2003modeling,
netzer2008mm}, но адаптирован к специфике онлайн-экспериментов и
краткосрочных интервенций.

\subsection*{Выводы по обзору}

Существующие подходы предоставляют богатый инструментарий для оценки
долгосрочных эффектов, однако каждое из направлений вынуждено идти на
компромисс между интерпретируемостью, масштабируемостью и надёжностью.
Длительные эксперименты дают наиболее точную оценку, но дороги и подвержены
эффектам выживаемости. Суррогатные модели ускоряют принятие решений ценой
интерпретируемости и опираются на сложно проверяемое предположение
суррогатности. Параметрические модели поведения и анализ выживаемости
строятся на сильных допущениях и плохо переносятся на разнородные
эксперименты. Uplift-моделирование требует значительных объёмов данных и
оценивает эффект на индивидуальном уровне, что затрудняет его применение
для стратегических решений.

Метод Orbitals закрывает обозначенный разрыв: он соединяет прозрачность
когортного анализа с предсказательной силой машинного обучения и сводит
оценку долгосрочного эффекта к анализу переходов между целеориентированными
поведенческими сегментами. Такая постановка естественно опирается на
операциональное определение долгосрочного эффекта и требования к методу,
которые сформулированы в главе <<Постановка задачи>>.